This work presents a conservative, modular post-segmentation pipeline for
automated split-error detection in connectomics data. The central research
question was: \emph{can a geometry-only, rule-based pipeline achieve high recall
on neural split detection without introducing excessive merge errors?}

On the CREMI Drosophila EM benchmark, the answer is clearly affirmative:
Precision = 1.000, Recall = 0.909, F1 = 0.952, with zero false positives. On
the XPRESS mouse white matter XNH challenge, a harder domain where the
segmented structures are elongated tubular axons with interior splits, the
best configuration (Experiment~24) achieves Recall = 0.9958 on the training
ground-truth set and Recall = 0.9879 on the fully held-out validation volume,
confirming generalization.

\paragraph{Key findings.}
Two innovations are responsible for the majority of the XPRESS recall improvement:

\begin{enumerate}
  \item The \textbf{skeleton-node graph} raises ground-truth coverage by 56.8
        percentage points over an endpoint-only graph in a single change, by
        enabling detection of interior splits along long axons.

  \item \textbf{Distance-conditioned scoring} prevents the proximity decay from
        suppressing long-range candidates whose alignment and continuity are
        strong, enabling the pipeline to accept splits at gaps up to 2000~nm.
\end{enumerate}

\paragraph{Challenges.}
The dominant open challenge is precision: Precision $\approx 0.003$ on XPRESS,
meaning approximately 312 false-positive merge proposals for every true positive.
The gap-stratified analysis in Section~\ref{sec:gap-stratified} localizes
this problem: the $[1500, 2000)$~nm bin alone contains 60.7\% of all false
positives while contributing only 16.3\% of true positives, and the
$[1000, 2000)$~nm long-range regime as a whole accounts for 74.8\% of
false positives at 25.6\% of true positives.  This stems from the geometric
feature space being insufficient to discriminate same-axon from cross-axon
long-range pairs at the density of myelinated white matter.  A machine-learned
filter (Experiment~20) improved precision $10\times$ (to 0.037) at the cost
of recall, but could not maintain recall $\geq 0.99$.  A neighborhood-density
re-ranker (Section~\ref{sec:density-reranker}) confirms that local crowding
is a real discriminator --- delivering up to 49\% relative precision gain at
recall 0.85 on both training and validation volumes --- but the effect
decays to $<$6\% at recall 0.99, showing that no single geometric signal
breaks the precision ceiling in this regime.  A per-bin ceiling analysis
(Section~\ref{sec:per-bin-threshold}) formalises this observation:
the precision ceiling achievable by any gap-aware threshold strategy
on the existing (composite, gap\_bin) feature pair is only $+$1.3\% at
recall 0.99 but rises to $+$104\% at recall 0.85 and $+$421\% at recall
0.70.  At recall 0.99 the hardest-to-recover TPs are distributed across
the same gap bins that hold the bulk of the FPs, so no per-bin
reallocation can remove FPs without sacrificing those TPs.  Meaningful
high-recall precision improvement therefore requires features
orthogonal to the current geometric scores (for example per-fragment
degree, as in Experiment~20) or access to raw image data.

The MinGapRule challenge was resolved in Experiment~24: disabling the rule
(min\_gap\_nm = 0) recovered all 37 short-gap false negatives and raised
training recall from 0.966 to 0.9958. The only remaining false negatives
(5 on training, 2 on validation) are caused by the CurvatureRule at gaps of
66--987~nm, where the local direction estimate is unreliable.

The coverage ceiling imposed by the 2000~nm search radius means at least 247
ground-truth pairs (with gap or PCA-tip distance exceeding this threshold)
cannot be reached by the current graph construction. The best achieved coverage
is 78.7\% on the training volume.

\paragraph{Scope and limitations.}
The pipeline's geometric scoring --- alignment, continuity, and the
composite's emphasis on local tangent agreement --- assumes that split
errors occur along structures with a consistent local axis. This assumption
holds strongly for myelinated axons in white matter (the XPRESS domain) but
weakens in two regimes. First, in higher-resolution EM of cortical gray
matter, splits frequently occur at approximately orthogonal angles --- for
example, orphaned dendritic spines whose connecting neck meets the parent
dendrite at a near-right angle. The DirectionReversalRule and CurvatureRule
would reject many such true merges. Second, the pipeline does not currently
model merge errors (two distinct neurons fused into one segment); it is
deliberately scoped to the split-correction half of the proofreading
problem. Extending to orthogonal-break morphologies would require
introducing learned morphological priors or conditioning the scoring on
estimated structure type (axon vs dendrite vs spine); extending to merge
detection would require a cut-proposal module downstream of fragment
extraction. Both are natural follow-ons and are discussed in
Chapter~\ref{ch:futurework}.

\paragraph{Summary of contributions.}
The pipeline is fully implemented, tested (428 tests, 0 failures), documented,
and open-source. It includes a YAML configuration system, six processing stages,
multi-format export, a ground-truth evaluation framework for two benchmarks,
an ML-based post-filter, and a visual diagnostic report generator with a
ground-truth TP sampling category that closes the scientific validation loop by
confirming that accepted merges correspond to anatomically real splits.
Twenty-eight experiments with detailed quantitative logs trace the development
of each component, including a gap-stratified characterization of the precision
challenge (Section~\ref{sec:gap-stratified}), a neighborhood-density re-ranker
that establishes an empirical bound on single-feature geometric precision
recovery (Section~\ref{sec:density-reranker}), and a per-bin ceiling analysis
that sets the theoretical upper bound for any gap-aware threshold strategy
(Section~\ref{sec:per-bin-threshold}).
